\documentclass[11pt,a4paper,twoside]{article}
\usepackage[utf8]{inputenc}
\usepackage[english]{babel}
\usepackage{actuarialangle, actuarialsymbol}
\usepackage{amsmath,amsfonts,amsthm,amssymb,mathtools}
\begin{document}

State the formula's for commutation figures and some important identities (but not specific insurance products)


For each age \( x \in \{ 0, 1, \ldots, w \} \), the commutation figures for an insurance with cause of elimination death are defined by

\(D_x\) := \(l_x v^x\) discounted number of the living
\(C_x\) := \(d_x v^{x+1} \) discounted number of the deaths
\(N_x\) := \(\sum_{k = x}^w D_k\) cumulated discounted number of the living
\(M_x\) := \(\sum_{k = x}^w C_k\) cumulated discounted number of the deaths
\(S_x\) := \(\sum_{k = x}^w N_k\) sum of second order of discounted number of the living
\(R_x\) := \(\sum_{k = x}^w M_k\) sum of second order of discounted number of the deaths

\( C_x = vD_x - D_{x+1} \)
\( M_x = D_x - dN_x \)
\( R_x = N_x - d S_x \)

The recursive computation (useful for computers/going fast in the exam)
Discounted number of the living:
\(x = 0 : \quad D_0 = l_0\)
\(x > 0 : \quad D_x = v^x\, l_x = v\, v^{x-1}\, l_{x-1}\, p_{x-1} = v\, D_{x-1}\, p_{x-1}\)

Cumulated discounted number of the living:
\(x = w : \quad N_w = D_w\)
\(x < w : \quad N_x = N_{x+1} + D_x\)

Sum of second order of the number of the living:
\(x = w : \quad S_w = N_w \)
\(x < w : \quad S_x = S_{x+1} + N_x\)



State the premia for specific life insurance products in terms of the commutation figures

Insurance with duration \( n \) and an insured sum of \( 1 \) , \( n \leq \omega + 1 - x \), in advance \\
\[ \Ex[n]{x} = \mathbb{E}[v^{n} \mathbf{1}_{(n, \infty)}(T_x)] = v^n \px[n]{x} = v^n \frac{l_{x+n}}{l_x} = \frac{D_{x+n}}{D_x} \]

Annual temporary annuity for life (for a living person) to be paid in advance with duration
\( n \leq \omega + 1 - x \) and an annual payment of 1
\[ \ax**{\endowxn} = \mathbb{E}[\sum_{k=1}^n \Ex[k]{x}] = \frac{1}{D_x}(N_{x+1} - N_{x+n+1}) \]
Annual lifelong annuity of \( 1 \) to be paid in arrear with duration \( n = \omega + x - 1 \); last
payment at the end of year of death:
\[ a_x = \frac{N_{x+1}}{D_x} = \frac{1}{D_x}(N_x - D_x) = \ax**{x} - 1\]

Yearly temporary annuity for life of \( 1 \) in advance, deferred by \( m \leq n \) years, with 
duration \( n - m \) years, \( n \leq \omega + 1 - x \) 
\begin{align*}
  \ax**[m|]{x:\angl{n - m}} &= \mathbb{E}[\sum_{k=m}^{n-1} v^k \mathbf{1}_{(k, \infty)}(T_x)]  \\
 &= \ax**{\endowxn} - \ax**{x:\angl m} = \frac{1}{D_x}(N_{x+m} - N_{x+n})
\end{align*}

Yearly lifelong annuity of \( 1 \) for life, paid in advance, deferred by \( m \leq \omega + 1 - x \) years,
with \( n = \omega + 1 - x \)
\[ \ax**[m|]{x} = \frac{N_{x+m}}{D_x} \]

Yearly temporary annuity for life of duration \( n \leq \omega + 1 - x \), to be paid in advance, with
an increasing payment beginning with \( 1 \) and a yearly increase of \( 1 \):
\begin{align*}
  (\Ia**)_{x:\angl n} &= \sum_{k=0}^{n-1} \ax**[k]{x:\angl{n - k}} = \frac{1}{D_x}(\sum_{k=0}^{n-1} N_x - N_{x+n}) \\
                      &= \frac{1}{D_x}(S_x - S_{x+n} - n N_{x+n})
\end{align*}

Yearly lifelong annuity for life, to be paid in advance, with a guaranteed period of \( g \) years
and a payment of \( 1 \):
\[ \ax**[][g]{x} = \ax**{g} + \ax[g|]{x} = \frac{1 - v^g}{1 - v} + \frac{N_{x+g}}{D_x} \]

\( n \)-year term life insurance with insured sum \( 1 \), where the benefits are due at the end of
the year of death, \( n \leq \omega + 1 - x \):
\begin{align*} 
  \Ax[n]{x} &= \mathbb{E}[\sum_{k=0}^{n-1} v^{k+1} \mathbf{1}_{(k, k+1]}(T_x)] = \sum_{k=0}^{n-1} v^{k+1} \px[k]{x} \qx{x+k} \\
 &= \frac{1}{D_x} \sum_{k=0}^{n-1} C_{x+k} = \frac{1}{D_x}(M_x - M_{x+n}) 
\end{align*}

Whole life insurance with insured sum \( 1 \) with \( n = \omega + 1 - x \):
\[ \Ax = \frac{M_x}{D_x} = 1 - d \frac{N_x}{D_x} \]

Term life insurance deferred by \( m \leq n \) years with duration \( n - m \) and insured sum \( 1 \),
due at the end of the year of death, \( n \leq \omega + 1 - x \):
\[ \Ax[m|n-m]{x} = \Ax[n]{x} - \Ax[m]{x} = \frac{1}{D_x}(M_{x+m} - M_{x+n}) \]

Whole life insurance deferred by \( m \leq \omega + 1 - x \) years with \( n = \omega + 1 - x \):
\[ \Ax[m]{x} = \frac{M_{x+m}}{D_x} \]

Term life insurance for \( n \) years with an increasing insured sum of \( 1 \) in the first year and
an increase of \( 1 \) per year:
\[ (\IA)_{x:\angl n} = \sum_{k=0}^{n-1} \Ax[k|]{x} = \frac{1}{D_x}(R_x - R_{x+n} - n M_{x+n}) \]

An endowment insurance for a period \( n \leq \omega + 1 - x \) covers both the risks of life and
death.\\
In case the insured person survives, the insured sum is paid out at the end of the term \( n \);\\
if death occurs, then it is paid immediately.
\[ \Ax{x:\angl n} = \Ex[n]{x} + \Ax[n]{x} = \frac{1}{D_x}(D_{x+n} + M_x - M_{x+n}) = 1 - d \ax**{\endowxn} \]


State recursive formula's for some life insurance products

Payment of the assured sum at the end of year of death, discretized:
\[ \Ax{x} = v\qx{x} + v\px{x} \Ax{x+1} \]

Immediate payment of the insured sum in the moment of death, continuous:
\[ \Ax*{x} = \mathbb{E}[v^{T_x}| T_x \leq h] \qx[h]{x} + v^h \Ax*{x+h} \px[h]{x} \]

Annuity with yearly payment of \( 1 \), payable in advance:
\[ \ax**{x} = 1 + v \ax**{x+1} \px{x} \]


List methods for recurring premium payments, instead of single ones

Recurring premia have the feature, that the policy holder pays premiums (either constant or dynamic) 
at certain points in time (monthly, annually, ...) until
1. death (for example private health insurance) \\
2. a contractually agreed upon point in time in case of survival, otherwise until death \\

Recurring premiums can be modeled as an annuity for life which is to be paid to the 
insurance company. Therefore we have for the present value of the premia: \\
\( \premium{x} \ax**{x} \) (lifelong premium payments) \\
\( \premium{x} \ax**{x:\angl m} \) (abbreviated premium payment period, m years) \\
Let \( \mathbb{E}[B_x^L] \) be, as before, the expected present value of the insurance benefits. \\
Then it follows: \\
\( \frac{\mathbb{E}[B^L_x]}{\ax**{x}} \) (lifelong premium period), \\
\( \frac{\mathbb{E}[B^L_x]}{\ax**{x: \angl m}} \) (lifelong premium period), \\

Consider a more involved example of a dynamic premium, say \( P_x^k \), \( k \in \{1, \dots, m\} \), payable in advance, 
increase yearly by a factor of \( c \% \). \\
\begin{align*}
  \sum_{k=0}^{m-1} v^k P_x^{k+1} \mathbb{P}[T_x > k] &= \sum_{k=0}^{m-1} v^k (1 + \frac{c}{100})^k P^1_x \mathbb{P}[T_x > k]  \\
  &= P^1_x \ax**{x:\angl m}
\end{align*}
where \( \ax**{x:\angl m} \) uses the modified factor discount factor \( v(1 + \frac{c}{100}) \).

Typically, insurances can be paid for not only annually, but also monthly, 
per quarter or half-yearly, in general: \( k \) premiums per year. \\
If the annual premium \( P_x \) is supposed to be paid for by \( k \) premiums in the amount of \( \frac{1}{k}P_x \)
per year, payable in advance, then the calculation of premium involves 
\[ \ax**{x:\angl m}[(k)] = \mathbb{E}\left[\sum^{mk - 1}_{j=0} v^{\frac{j}{k}} \frac{1}{k} \mathbf{1}_{(\frac{j}{k}, \infty)}(T_x) \right] \]
for an annuity for life with \( k \) payments per year, in advance, to be paid for \( m \) years in the
amount of \( \frac{1}{k} \).

Therefore the principle of equivalence implies for \\
\( P_x \ax**{x:\angl m}[(k)] = \mathbb{E}[B^L_x] \) or \( P_x \ax**{x}[(k)] = \mathbb{E}[B^L_x] \) \\
with \( k \)-yearly premium payments, is simply 
\( P_x = \frac{\mathbb{E}[B^L_x]}{\ax**{x:\angl m}[(k)]} \). \\
These premiums \( \frac{1}{k}P_x \), due \( k \) times per year, are called real inter-annual premiums. \\
In practice, inter-annual premiums are often calculated 
approximatively on the basis of annual premiums also then referred to as false inter-annual premiums: \\
Since the whole annual premium reaches the insurance not at the beginning of the year,
but peu a peu during the whole year, a surcharge for installments is raised on the
premiums \( \frac{1}{k}P_x \). This method is criticized as it is an implicit loan! \\

Some notational remarks: \\
\( m \) abbreviated premium payment period \\
\( P_{x,m}(\mathbb{E}[B^L_x]) \) annual net premium for an \( x \)-year old that is to be paid
for \( m \) years for a product with expected amount
of benefits \( \mathbb{E}[B_x^L] \).


Give some examples of standard life insurance products with recurring payments

Pure endowment insurance (benefits only in case of survival) for \( n \) years with insured sum \( 1 \)
\( m \leq n < \omega - x \), payable in advance:
\[ P_{x,m}(\Ex[n]{x}) = \frac{\Ex[n]{x}}{\ax**{x:\angl m}} = \frac{D_{x+n}}{N_x - N_{x+m}} \]

Temporary annuity for life, to be paid out in advance, deferred by \( m \leq n < \omega - x \) years
with a duration of \( n - m \) and a yearly annuity of \( 1 \), premium payments during the deferral period:
\[ P_{x,m}(\ax**[m|]{x:\angl m}) = \frac{\ax**[m|]{x:\angl{n - m}}}{\ax**{x:\angl m}} = \frac{N_{x+m} - N_{x+n}}{N_x - N_{x+m}}  \]

Annuity for life, to be paid out in advance, deferred by \( m < \omega - x \) years and with a yearly
annuity of \( 1 \), premium payments during the deferral period:
\[ P_{x,m}(\ax**[m|]{x}) = \frac{\ax**[m|]{x}}{\ax**{x:\angl m}} = \frac{N_{x+m}}{N_x - N_{x+m}}  \]

Insurance on a fixed date (Term-fix Insurance) with duration \( n \), insured sum \( 1 \) and
premium payment period \( m \leq n \):
\[ P_{x,m}(v^n) = \frac{v^n}{\ax**{x:\angl m}} = v^n \frac{D_x}{N_x - N_{x+m}} \]
Note that only the number of premiums is unknown. \\

(Temporary) term life insurance for \( n \) years with insured sum \( 1 \), benefits are to be paid
at the end of the year of death, premium payment period \( m \leq n \):
\[ P_{x,m}(\Ax[n]{x}) = \frac{\Ax[n]{x}}{\ax**{x: \angl m}} = \frac{M_x - M_{x+n}}{N_x - N_{x+m}} \]



Define the different kinds of cost encountered in insurance contracts and how to include them in the premia

We distinguish the following costs:
\begin{enumerate}
  \item \( \alpha^Z, \alpha^\gamma \)-costs, : direct and indirect acquisition costs respectively, direct depend on the conclusion of a contract and are due immediately (commissions, costs for doctors, etc.) and indirect (salaries and costs for office space for the sales force)
  \item \( \gamma \)-costs : administrative costs, office rent, IT licenses
  \begin{enumerate}
    \item sometimes there are split into \( \gamma_1, \gamma_2 \) for concurrent with premium payments and costs afterwards.
  \end{enumerate}
\item \( \beta \)-costs : collection costs, incurred from money transactions (Bank fees, Employees in the collection department)
\end{enumerate}
The greek letters usually represent the percentage of some calculation base sum of premiums, sum assured or annuity,
which are denoted by \( GA^Z, GA^\gamma, GG \) for \( \alpha^Z, \alpha^\gamma, \gamma \) respectively. \\

In the calculation of premiums, costs of first order are used. The real costs are called
costs of second order, in analogy to the other calculation bases.
The resulting premia with the surcharges is also referred to as sufficient premium. \\

Instead of handing these costs over to the policy holder as a single premium, these
acquisition costs are (partly) paid for in the form of surcharges to the current premiums,
so Acquisition costs are earned over time. \\

In the simplest setting, only using the direct acquisition costs this is also called Zillmer premium,
which must satisfy
\[ P^Z_x \ax**{x: \angl m} = \mathbb{E}[B^L_x] + \alpha^Z GA^Z  \]
typically \( GA^Z = m P_x^Z \) (depends on the sum of the premiums) so the premia would be given by
\[ P^Z_x = \frac{\mathbb{E}[B^L_x]}{\ax**{x: \angl m} - \alpha^Z m} \]

More generally, assume constant premia, 
due yearly in advance, 
policy term is \( n < \omega - x\); 
the premium payment period is \( m \leq n \) years:
\[ 
\premium[][a]{x} = \frac{1}{(1 - \beta) \ax**{x: \angl m}}( \mathbb{E}[B^L_x] + \alpha^Z GA^Z + \alpha^\gamma GA^\gamma + \gamma GG \ax**{x: \angl n}) 
\]
If the \( \gamma \)-costs are split, and the direct acquisition costs depend through broker fees on the overall premia this is instead 
\[ 
\premium[][a]{x} = \frac{1}{(1 - \beta) \ax**{x: \angl m} - \alpha^Z m}( 
\mathbb{E}[B^L_x] + \alpha^Z GA^Z + \alpha^\gamma GA^\gamma + GG (\gamma_1 \cdot \ax**{x: \angl m} + \gamma_2 \cdot \ax**[m|]{x: \angl{n - m}}) )
\]
if the insurance immediately only terminates with the death of the person, the \( \gamma_2 \) costs continue indefinitely 
\[ 
\premium[][a]{x} = \frac{1}{(1 - \beta) \ax**{x: \angl m} - \alpha^Z m}( 
\mathbb{E}[B^L_x] + \alpha^Z GA^Z + \alpha^\gamma GA^\gamma + GG (\gamma_1 \cdot \ax**{x: \angl m} + \gamma_2 \cdot \ax**[m|]{x}))
\]

(can also be derived from equivalence principle).


Elaborate on sum discounts and surcharges in the context of tariff premia

The tariff premium is the sufficient premium adjusted for surcharges and discounts. The
following components are to be considered:
Unit cost surcharge: On top of the already considered costs, there may be additional
costs which are fixed for every single contract (independent of insured sum and premium).
These can be considered by adding a fixed amount to the premium.

Surcharges or discounts for installments: If the sufficient annual premium is calculated for
yearly payments in advance, but the premium is paid during the year (e. g., monthly),
then a surcharge for installments is applied (artificial inter-annual premium).
If the sufficient annual premium is calculated for yearly payments in arrears, but the
premium is paid for over time during the year (e. g., monthly) before it is actually due,
then a discount for installments can be applied.

From an economic point of view, it therefore makes sense (and is objectively correct) to
reduce costs for high insured sums and increase costs for small insured sums

Because of the competition between insurance companies, also risk discounts have an
increased importance, e. g., a non-smoker tariff for a term life insurance
These are called risk discounts.

If we add insurance taxes on the tariff premium, we get the gross premium.
In Germany, there are no taxes on insurance premiums in life insurance, therefore tariff
premium and gross premium can be used synonymously.

On a last note, the German VAG requires the tariffication may not be changed according to high demand,
any difference in price or benefit must be analytically justifiable, statistically verifiable and
transparent for a third person.

Therefore the portfolio (of insurance contracts) is split into
disjoint subsets each using a different identical calculation base.
\end{document}
